\chapter*{前言}
\phantomsection\label{front:preface}
\addcontentsline{toc}{chapter}{前言}
\begin{center}\BookEpigraph{preface}\end{center}

一次请求没有收到答复，往往只需要再试一次。麻烦在于，没收到答复的人，并不知道上一次已经走到了哪里。服务可能尚未开始，也可能已经改变了账户、发出了消息，只是来不及把结果送回来。重试让系统能够继续工作，也把一个原本局部的问题带进了别人的账本。

幂等性常常以一个简短的要求出现：同一操作重复执行，不应产生额外效果。这句话适合写在需求里，却不足以直接指导实现。什么算同一操作，哪些变化属于它的效果，谁保存已经发生过的证据，证据丢失以后还能承诺什么，都需要在具体系统中说明。数据库里的唯一约束能挡住一些重复写入，却不能替另一个服务撤回已经完成的动作；一份成功记录能回答今天的重试，也未必经得起明天的备份恢复。

本书写给已经接触过接口、事务和消息队列，希望把这些疑问落实到代码中的工程师。我们从一个可以正常运行的小系统开始，让请求重复，让响应丢失，让进程停在不方便的位置，再查看留下的记录。前三部分让实现随着问题逐步改变，第四部分从总体架构出发，讲解服务、数据、复制、事务、执行运行时与多地域部署怎样组成完整系统。每增加一层机制，都要说明它依赖什么、保护什么，以及在哪一种执行顺序下仍会失效。

分析一次故障，需要把已经发生的动作和各方掌握的信息分别列清。数据库可能已经提交，调用方却仍在等待答复；进程恢复后，内存中的执行进度已经丢失，持久化记录还在。书中的实验会保留这些时刻的执行记录，方便读者对照代码，检查重试与恢复的判断依据。

一项工作经由多人或多个进程接力完成时，后续处理总要依赖前面留下的信息。一次确认没有送到，或者一份记录没有及时更新，都可能使已经完成的事情重新进入待办列表。要让工作继续，又避免无意中重做，就得使接手者有办法核实先前的进度，并明确证据不足时该如何处置。

